\documentclass[10pt,a4paper]{amsart}
\usepackage[margin=19mm]{geometry}
\usepackage{amsmath,amssymb,mathtools}
\usepackage[scr=rsfs,cal=euler]{mathalfa}
\usepackage{xfrac}
\usepackage[unicode,hidelinks,bookmarks=false]{hyperref}
\newcommand{\ie}{i.e.\ }
\newcommand{\cf}{cf.\ }
\newcommand{\resp}{respectively\ }
\newcommand{\Subsection}{\subsection}
\providecommand{\nobreakdash}{\nobreak\hbox{-}}
\setlength{\emergencystretch}{2em}
\multlinegap0pt
\newtheorem{comment}{Comment}
\usepackage{amssymb}
\usepackage{xcolor}
\newtheorem{proposition}{Proposition}
\newtheorem{lemma}{Lemma}
\def\A{{\mathfrak A}\, }
\title{$*$-Jordan-type maps on alternative $*$-algebras}
\author{Aline J. O. Andrade, Bruno L. M. Ferreira, Liudmila Sabinina}
\date{Source: arXiv:2307.00002v1 (CC BY 4.0)}
\begin{document}
\maketitle

\noindent\emph{Source and licence.} $*$-Jordan-type maps on alternative $*$-algebras. Version arXiv:2307.00002v1 (CC BY 4.0), distributed by arXiv under the Creative Commons Attribution 4.0 International licence, http://creativecommons.org/licenses/by/4.0/, as recorded in the arXiv licence metadata for this exact version and shown on the abstract page https://arxiv.org/abs/2307.00002v1. Full licence notice: \texttt{licenses/arxiv-2307.00002-CC-BY-4.0.txt}.\par\smallskip

\noindent\textit{Editorial scope.} Counted unit: the article's lemma lema2 (source0.tex lines 263--265). The article's complete original proof of it is delivered with it (source0.tex lines 267--328, one \texttt{proof} environment of 62 lines). No mathematics is written, repaired or re-derived: the statement and the proof are the article's own lines, in the article's own notation, and no retained line is substituted or re-flowed.  Retained uncounted as the source-local context the proof needs: lines 156--175; lines 245--247.  Nothing was omitted from any retained span, so the package carries no omission record.  No retained line contains a citation, so the wrapper carries no bibliography and no bibliography entry.  No retained line contains a figure environment, an image-inclusion command or drawing code, so no image is delivered and the package declares no asset dependency.  Editorial formatting only, in addition: an AMS \texttt{amsart} preamble that restates the article's own declarations and macros used by the retained lines, namely the environments \texttt{proposition}, \texttt{lemma} on separate counters, together with the standard packages those lines need.  The retained environments print in this document as Proposition 1, Lemma 1 in order of appearance, whereas the article prints the same items with the numbering of its own full document.  The complete native source of the article as distributed in the arXiv e-print for this exact version is bundled unchanged as \texttt{sources/143661/source0.tex}.
\begin{proposition}\label{Claim2} Let $\A$ and $\A'$ be two alternative $*$-algebras and $\varphi: \A \rightarrow \A'$ a bijective map which satisfies
\[
\varphi(q_{n*}(\xi, \ldots,\xi, a,b)) = q_{n*}(\varphi(\xi),\ldots,\varphi(\xi),\varphi(a),\varphi(b))
\]
for all $a, b \in \A$ and $\xi \in \left\{1_{\A}, e_1, e_2\right\}$. Let $x,y$ and $h$ be in $\A$ such that $\varphi(h) = \varphi(x) + \varphi(y)$. Then,
 given $z \in \A$,
$$
\begin{aligned}
\varphi(q_{n*}(t,...,t,h,z)) = \varphi(q_{n*}(t,...,t,x,z)) + \varphi(q_{n*}(t,...,t,y,z))
\end{aligned}
$$
and
$$
\begin{aligned}
\varphi(q_{n*}(t,...,t,z,h)) = \varphi(q_{n*}(t,...,t,z,x)) + \varphi(q_{n*}(t,...,t,z,y))
\end{aligned}
$$
for $t = 1_\A$ or $t = e_i$, $i = 1,2$.

\end{proposition}

\begin{lemma}\label{lema1} For any $a_{11} \in \A_{11}$ and $b_{22} \in \A_{22}$, we have 
$$\varphi(a_{11} + b_{22}) = \varphi(a_{11}) + \varphi(b_{22}).$$
\end{lemma}

\begin{lemma}\label{lema2}
For any $a_{12} \in \A_{12}$ and $b_{21} \in \A_{21}$, we have $\varphi(a_{12} + b_{21}) = \varphi(a_{12}) + \varphi(b_{21})$.
\end{lemma}

\begin{proof}
Since $\varphi$ is surjective, given $\varphi(a_{12})+\varphi(b_{21}) \in \A'$ there exists $t \in \A$ such that 
$\varphi(t) = \varphi(a_{12})+\varphi(b_{21})$, with $t=t_{11}+t_{12}+t_{21}+t_{22}$. Now, by Proposition \ref{Claim2}
$$
\begin{aligned}
\varphi\left(q_{n*}\left(e_1,...,e_1,\frac{1}{2^{n-1}}e_1,t\right)\right) &= \varphi\left(q_{n*}\left(e_1,...,e_1,\frac{1}{2^{n-1}}e_1,a_{12}\right)\right) \\
& +\varphi\left(q_{n*}\left(e_1,...,e_1,\frac{1}{2^{n-1}}e_1,b_{21}\right)\right) \\ 
&= \varphi(e_1a_{12} + a_{12}e_1) + \varphi(e_1b_{21} + b_{21}e_1) \\
&= \varphi(a_{12}) + \varphi(b_{21}) = \varphi(t).
\end{aligned}
$$
Since $\varphi$ is injective, $e_1t + te_1 = t,$ that is,
$2t_{11} + t_{12} + t_{21} = t_{11}+t_{12}+t_{21}+t_{22}.
$ Then $t_{11} = t_{22} = 0$.
%Now, observe that, for $c_{12} \in \A_{12}$, $q_{n*}(e_1,...,e_1,c_{12},a_{12}) \in \A_{11} + \A_{21}$ and $q_{n*}(e_1,...,e_1,c_{12},b_{21}) \in \A_{11} + \A_{12}$. Then,  $q_{n*}(e_1, \cdots, e_1, a_{12},c_{12})=2^{n-2}(c_{12}a_{12}+a_{12}c_{12}^*) \in \A_{11}+\A_{21},$ $q_{n*}(e_1, \cdots,e_1,c_{12},b_{21})\in \A_{11}+\A_{12}$ and  






Now, observe that, for $c_{12}\in \A_{12},$ $q_{n*}\left(e_2, \ldots, e_2,\dfrac{1}{2^{n-2}}e_2,t,c_{12}\right) = t_{21}c_{12}+t_{12}c_{12}+c_{12}t_{21}^*+c_{12}t_{12}^*.$ Thus, $q_{n*}\left(e_2,\ldots,e_2,\dfrac{1}{2^{n-2}}e_2,a_{12},c_{12}\right) = a_{12}c_{12}+c_{12}a_{12}^* \in \A_{11}+\A_{21}$ and $q_{n*}\left(e_2, \ldots,e_2, \dfrac{1}{2^{n-2}}e_2,b_{21},c_{12}\right)=b_{21}c_{12}+c_{12}b_{21}^* \in \A_{21}+\A_{22}.$ Moreover, 
\[q_{n*}\left(e_2, \ldots, e_2, \dfrac{1}{2^{n-2}}e_2, q_{n*}\left(e_2, \ldots, e_2,\dfrac{1}{2^{n-2}}e_2,t,c_{12}\right),e_2\right) =2(t_{21}c_{12}+(t_{21}c_{12})^*),
\]
thus 
\begin{align*}
\varphi& \left(q_{n*}\left(e_2,\ldots,e_2, \dfrac{1}{2^{n-2}}e_2,q_{n*}\left(e_2, \ldots,e_2, \dfrac{1}{2^{n-2}}e_2,t,c_{12}\right),e_2\right)\right) \\
=& \ \varphi\left(q_{n*}\left(e_2,\ldots,e_2,\dfrac{1}{2^{n-2}}e_2,q_{n*}\left(e_2, \ldots,e_2,\dfrac{1}{2^{n-2}}e_2,a_{12},c_{12}\right),e_2\right)\right)\\
& \ +\varphi\left(q_{n*}\left(e_2,\ldots,e_2,\dfrac{1}{2^{n-2}}e_2,q_{n*}\left(e_2, \ldots,e_2,\dfrac{1}{2^{n-2}}e_2,b_{21},c_{12}\right),e_2\right)\right) \\
=& \ \varphi(0) + \varphi(2(b_{21}c_{12}+(b_{21}c_{12})^*))=\varphi(2(b_{21}c_{12}+(b_{21}c_{12})^*)).
\end{align*}
Therefore, $(t_{21}-b_{21})c_{12}+c_{12}^*(t_{21}^*-b_{21}^*)=0,$ for all $c_{12}\in \A_{12}.$ Consider $ic_{12}\in \A_{12},$ we have $(t_{21}-b_{21})c_{12}-c_{12}^*(t_{21}^*-b_{21}^*)=0.$ Thus, $(t_{21}-b_{21})c_{12}=0$ implies $(t_{21}-b_{21})(\A e_2) = 0.$ Using ($\spadesuit$), we obtain $t_{21}-b_{21}=0,$ that is, $t_{21}=b_{21}.$ Note that, with similar calculations, if we replace $c_{12}$ by $c_{21}$ and $e_1$ by $e_2,$ obtain that $t_{12} = a_{12}.$
\begin{comment}
Then, by Proposition \ref{Claim2} and Lemma \ref{lema1}, we obtain
$$
\begin{aligned}
\varphi(q_{n*}(e_1,...,e_1,t,c_{12})) &= \varphi(q_{n*}(e_1,...,e_1,a_{12},c_{12})) + \varphi(q_{n*}(e_1,...,e_1,b_{21},c_{12})) \\
&= \varphi(q_{n*}(e_1,...,e_1,a_{12},c_{12})+q_{n*}(e_1,...,e_1,b_{21},c_{12})).
\end{aligned}
$$



 



By injectivity of $\varphi$ we have
$$
q_{n*}(e_1,...,e_1,t,c_{12})=q_{n*}(e_1,...,e_1,a_{12},c_{12})+q_{n*}(e_1,...,e_1,b_{21},c_{12}),
$$
that is,
$$
t_{21}c_{12}+c_{12}t_{12}^* = b_{21}c_{12}+c_{12}a_{12}^*.
$$
Therefore,
$$
(t_{21} - b_{21})c_{12} = 0 \mbox{\,\, and \,\,} c_{12}(t_{12}^*-a_{12}^*) = 0.
$$
Finally, by $\left(\spadesuit\right)$ we conclude that $t_{12}=a_{12}$ and $t_{21}=b_{21}$.
\end{comment}
\end{proof}

\end{document}
